\section{Discussion}
\label{sec:discussion}

\subsection{What the archetypes mean, and where they break}

The six outfield archetypes describe three axes rather than six independent types, since
standardising within position group before a two-way split forces the two centroids in
each group to be mirror images. Each axis orders players from low to high attacking
involvement within their listed group, and in every group the high involvement end is
distinguished by crossing and chance creation while the low involvement end is
distinguished by defensive actions or, in the case of defenders, by the absence of anything
positive at all. That an ordering exists is a real result. That it is a single ordering
rather than a taxonomy is the limit the data now imposes.

The evidence on how far to trust these divisions is unusually consistent. Bootstrap
resampling, replication on an independent season and shrinkage toward the position mean are
three different tests of three different failure modes, and they agree on which partitions
hold. The midfield split survives all three. The defender and forward splits fail all
three, and the forward split reaches an adjusted Rand index below zero in the least
favourable bootstrap resample, meaning the partition recovered there was worse than a
random relabelling. A reader who wants one sentence of practical guidance should take this
one: use the global attacking involvement axis, and do not characterise an individual
defender or forward by his within-group cluster.

The goalkeeper analysis deserves separate emphasis because it failed for a reason worth
recording. The partition that the selection rule returned explains more variance in team
points per match than in save percentage. In other words the goalkeeper clusters recover
the quality of the team in front of the goalkeeper rather than the goalkeeper himself. With
post-shot expected goals removed, almost every surviving goalkeeper statistic is a function
of how much and what quality of shooting the defence permits, so this is close to
inevitable. Reporting the partition without that check would have produced plausible
looking goalkeeper archetypes that were actually a league table in disguise.

\subsection{Listed positions and data-driven roles}

The clusters do not recover listed positions, and the interesting part is where they
disagree. The global division places every defender and a majority of midfielders on one
side, with the remaining midfielders and every forward on the other, so it cuts the midfield
in half and leaves the defender and forward groups intact. The players it reassigns are
recognisable: holding midfielders sit with the defenders, and attacking midfielders sit
with the forwards. That the same players are also the ones a supervised classifier places
in the wrong group is corroboration from a second method, although the enrichment over
chance is modest because the hybrid flag is broad.

The similarity search makes the limitation concrete in a way the aggregate statistics do
not. A holding midfielder's nearest statistical neighbours are creative attacking
midfielders, because the surviving features contain nothing that describes his actual job.
Tackles by pitch third, blocks, clearances, pressures and passing volume are all gone, so
the only distinctive things left in his profile are set piece deliveries and key passes,
and the metric duly finds the players who share those. An inverted full back returns centre
backs for the same reason. These are not failures of the method but measurements of what
the data no longer contains, and they are the clearest illustration in the paper of what
the withdrawal cost.

\subsection{Limitations}

Season aggregates discard context: a tackle count does not say where the tackles happened,
and an expected goals total does not distinguish six poor chances from two good ones.
Nothing here is possession adjusted, which systematically inflates the defensive counts of
players in teams that spend more time without the ball, and with team strength correlating
strongly with the first principal component that is a real confound rather than a
theoretical one. Per-90 rates remain noisy above the minutes threshold, which is why the
shrinkage analysis matters and why its within-group verdict should be believed. Merging a
second expected goals provider introduces a model whose assumptions differ from the one
that produced the original values, and although we never mix the two sources for the same
quantity, the creation block and the shooting block now come from different modelling
traditions. Roughly three percent of records failed to match across the two sources and
carry missing values for the Understat block. The study covers one league and two seasons.
Finally, the clusters are descriptive and carry no causal content: an archetype label
records how a player's season looked, not why, and certainly not how he would perform in a
different team.

\subsection{Applications}

For scouting and recruitment the honest use of this work is narrow but real. The global
involvement axis and the club level signatures are stable enough to support questions about
squad composition, such as whether a club's minutes are concentrated in one part of the
distribution, and the twenty panel comparison makes that legible at a glance. The similarity
search is usable for forwards and full backs, where the surviving features capture most of
what defines the role, and should not be used for central midfielders or centre backs,
where they do not. For tactical analysis the correlation divergence result is the most
directly applicable: any analysis that pools across positions when computing relationships
between statistics is averaging over sign reversals.

\subsection{Future work}

The obvious extension is data. Event and tracking data would restore everything the
withdrawal removed and more, at the cost of the reproducibility that a public source
provides. Within public data, match level statistics remain available in a reduced form and
would support the role drift analysis that was scoped out of this study, tracking whether a
player's profile moves between archetypes across a season rather than assuming a season is
homogeneous. Extending to other leagues would test whether the two-cluster result is a
property of the reduced feature set or of the Premier League specifically, and we expect the
former. A more careful treatment of possession adjustment would address the confound
between defensive volume and team strength directly. Finally, the empirical Bayes machinery
developed here is more broadly useful than the single use it is put to: any per-90 ranking
computed over partial seasons is subject to the same problem, and most published ones do
not correct for it.
